\begin{ouverture}[geometrie]{16}{Triangles et cercle circonscrit}{Partie 3 · Espace et géométrie}
{Constructions · Triangles particuliers · Somme des angles · Cercle circonscrit}
\objectif{Construire des triangles à partir de longueurs et d’angles.}
\objectif{Utiliser la somme des angles et les angles des triangles particuliers.}
\objectif{Construire le cercle circonscrit à un triangle.}
\prerequis{Compas ; rapporteur ; médiatrice d’un segment.}
\end{ouverture}
\pagestyle{otMain}\essentiel\fontsize{11}{15}\selectfont
\phantomsection\label{cours-16-1}
\essSection{1}{Construire et discuter les données}
\begin{savoirfaire}[Trois longueurs]Pour construire ABC avec $AB=6$ cm, $AC=4$ cm et $BC=5$ cm : tracer [AB] ; tracer un arc de centre A et de rayon 4 cm, puis un arc de centre B et de rayon 5 cm. Leur point d’intersection est C.\end{savoirfaire}
\label{fig-cours-16-1-cours-construction-triangle}\figureGeo{cours-construction-triangle}

\begin{propriete}Trois longueurs ne permettent pas toujours de construire un triangle : la plus grande doit être plus petite que la somme des deux autres. Avec 3 cm, 4 cm et 8 cm, les arcs ne se rencontrent pas, car $3+4<8$.\end{propriete}
\essSection{2}{Triangles particuliers et angles}
\begin{definition}Un triangle isocèle a deux côtés de même longueur : ses deux angles à la base ont la même mesure. Un triangle équilatéral a trois côtés de même longueur : ses trois angles mesurent $60^\circ$. Un triangle rectangle a un angle droit.\end{definition}
\clearpage
\phantomsection\label{cours-16-2}
\essSection{3}{Somme des angles d’un triangle}
\begin{propriete}Dans tout triangle, la somme des mesures des trois angles vaut $180^\circ$.\end{propriete}
\label{fig-cours-16-2-cours-somme-angles}\figureGeo{cours-somme-angles}

\begin{exemple}Si deux angles mesurent $48^\circ$ et $67^\circ$, le troisième mesure $180-48-67=65^\circ$. Dans un triangle rectangle, les deux angles aigus ont une somme de $90^\circ$.\end{exemple}
\essSection{4}{Médiatrices et cercle circonscrit}
\begin{propriete}Les trois médiatrices des côtés d’un triangle se coupent en un même point O. Ce point est à égale distance des trois sommets : c’est le centre du cercle qui passe par A, B et C, appelé cercle circonscrit au triangle.\end{propriete}
\label{fig-cours-16-2-cours-cercle-circonscrit}\figureGeo{cours-cercle-circonscrit}

\begin{savoirfaire}[Construire le cercle circonscrit]Deux médiatrices suffisent : leur point d’intersection est le centre O ; le rayon est OA.\end{savoirfaire}

\clearpage\rubrique{Découverte}
\begin{blocLivre}{ex-16-1}\begin{exercice}[1]{Le troisième angle}
Dans le triangle ABC, $\widehat{A}=48^\circ$ et $\widehat{B}=67^\circ$. Calculer $\widehat{C}$.
\label{fig-ex-16-1}\figureGeo{tri-1}
\end{exercice}\end{blocLivre}
\begin{blocLivre}{ex-16-2}\begin{exercice}[1]{Triangle isocèle}
Le triangle ABC est isocèle en A et $\widehat{A}=40^\circ$. Calculer les deux autres angles.
\label{fig-ex-16-2}\figureGeo{tri-2}
\end{exercice}\end{blocLivre}
\begin{blocLivre}{ex-16-3}\begin{exercice}[1]{Triangle équilatéral}
Un triangle équilatéral a des côtés de 4 cm. Donner la mesure de ses angles et son périmètre.
\label{fig-ex-16-3}\figureGeo{tri-3}
\end{exercice}\end{blocLivre}
\begin{blocLivre}{ex-16-4}\begin{exercice}[1]{Triangle rectangle}
Le triangle ABC est rectangle en C et $\widehat{A}=35^\circ$. Calculer $\widehat{B}$.
\label{fig-ex-16-4}\figureGeo{tri-4}
\end{exercice}\end{blocLivre}
\begin{blocLivre}{ex-16-5}\begin{exercice}[1]{Trois longueurs}
Construire le triangle ABC tel que $AB=6$ cm, $AC=4$ cm et $BC=5$ cm.
\label{fig-ex-16-5}\figureGeo{tri-5}
\end{exercice}\end{blocLivre}

\rubrique{Entraînement}
\begin{blocLivre}{ex-16-6}\begin{exercice}[2]{Possible ou pas ?}
Peut-on construire un triangle dont les côtés mesurent 3 cm, 4 cm et 8 cm ? Et 5 cm, 6 cm et 10 cm ?
\label{fig-ex-16-6}\figureGeo{tri-6}
\end{exercice}\end{blocLivre}
\begin{blocLivre}{ex-16-7}\begin{exercice}[2]{Deux côtés et un angle}
Construire le triangle ABC tel que $AB=5$ cm, $\widehat{BAC}=60^\circ$ et $AC=3$ cm.
\label{fig-ex-16-7}\figureGeo{tri-7}
\end{exercice}\end{blocLivre}
\begin{blocLivre}{ex-16-8}\begin{exercice}[2]{Un côté et deux angles}
Construire le triangle ABC tel que $BC=6$ cm, $\widehat{B}=50^\circ$ et $\widehat{C}=60^\circ$. Calculer $\widehat{A}$ sans le mesurer.
\label{fig-ex-16-8}\figureGeo{tri-8}
\end{exercice}\end{blocLivre}
\begin{blocLivre}{ex-16-9}\begin{exercice}[2]{Angles d’un isocèle}
Le triangle ABC est isocèle en A et ses angles à la base mesurent $72^\circ$. Calculer l’angle au sommet $\widehat{A}$.
\label{fig-ex-16-9}\figureGeo{tri-9}
\end{exercice}\end{blocLivre}
\begin{blocLivre}{ex-16-10}\begin{exercice}[2]{Construire le cercle circonscrit}
Reproduire le triangle ABC, puis construire son cercle circonscrit.
\label{fig-ex-16-10}\figureGeo{tri-10}
\end{exercice}\end{blocLivre}

\rubrique{Problèmes}
\begin{blocLivre}{ex-16-11}\begin{exercice}[3]{Trois villages}
Trois villages A, B et C veulent installer une antenne à la même distance des trois. Où la placer ?
\label{fig-ex-16-11}\figureGeo{tri-11}
\end{exercice}\end{blocLivre}
\begin{blocLivre}{ex-16-12}\begin{exercice}[3]{Le cercle d’un triangle rectangle}
Le triangle ABC est rectangle en A. Construire son cercle circonscrit. Que remarque-t-on sur la position du centre ?
\label{fig-ex-16-12}\figureGeo{tri-12}
\end{exercice}\end{blocLivre}
\begin{blocLivre}{ex-16-13}\begin{exercice}[3]{Isocèle et équilatéral}
Le triangle ABC est isocèle en A avec $\widehat{B}=65^\circ$. On construit, de l’autre côté de (AC), le triangle équilatéral ACD. Calculer $\widehat{BAC}$, puis $\widehat{BAD}$.
\label{fig-ex-16-13}\figureGeo{tri-13}
\end{exercice}\end{blocLivre}
\begin{blocLivre}{ex-16-14}\begin{exercice}[3]{Dans un triangle équilatéral}
ABC est un triangle équilatéral. La bissectrice de l’angle $\widehat{A}$ coupe le côté [BC] en H. Calculer $\widehat{BAH}$, puis $\widehat{AHB}$.
\label{fig-ex-16-14}\figureGeo{tri-14}
\end{exercice}\end{blocLivre}
\begin{blocLivre}{ex-16-15}\begin{exercice}[3]{Deux angles obtus ?}
Zoé veut construire un triangle ABC avec $\widehat{A}=100^\circ$ et $\widehat{B}=95^\circ$. Est-ce possible ?
\label{fig-ex-16-15}\figureGeo{tri-15}
\end{exercice}\end{blocLivre}

\clearpage
